% ============================================================
%  SOAL + PEMBAHASAN  --  FISIKA  --  SSU-ITB 2024
%  Tipe A (1-20) + Tipe B (21-39) + Soal Baru (40) | Paket 1
% ============================================================

\documentclass[11pt]{article}

\usepackage[utf8]{inputenc}
\usepackage[T1]{fontenc}
\usepackage[a4paper,margin=2.5cm,top=2.8cm]{geometry}
\usepackage{amsmath,amssymb}
\usepackage{graphicx}
\usepackage{enumitem}
\usepackage{multicol}
\usepackage{xcolor}
\usepackage[most]{tcolorbox}
\usepackage{fancyhdr}
\usepackage{booktabs}
\usepackage{icomma}
\usepackage{array}

% --- Warna Fisika ---
\definecolor{mapelcolor}{RGB}{40,90,150}
\definecolor{mapelbg}{RGB}{230,240,255}
\definecolor{pembcolor}{RGB}{100,50,130}
\definecolor{pembbg}{RGB}{245,238,252}
\definecolor{gold}{RGB}{210,160,30}

\pagestyle{fancy}\fancyhf{}
\lhead{\small SSU-ITB 2024 --- Pembahasan Fisika}
\rhead{\small Paket 1}
\cfoot{\small\thepage}
\renewcommand{\headrulewidth}{0.4pt}

\newtcolorbox{soal}[1]{
  breakable,
  colback=mapelbg, colframe=mapelcolor,
  coltitle=white, fonttitle=\bfseries\sffamily,
  title=Nomor #1,
  boxrule=0.8pt, arc=3pt,
  left=3mm, right=3mm, top=2mm, bottom=2mm}

\newtcolorbox{pembox}{
  breakable,
  colback=pembbg, colframe=pembcolor,
  coltitle=white, fonttitle=\bfseries\sffamily,
  title=Pembahasan,
  boxrule=0.8pt, arc=3pt,
  left=3mm, right=3mm, top=2mm, bottom=2mm}

\newcommand{\jawaban}[1]{%
  \noindent\colorbox{gold!30}{\parbox{\dimexpr\linewidth-2\fboxsep\relax}{%
    \textbf{Jawaban:} #1}}\par\smallskip}

\newcommand{\konsep}{\noindent\textit{Konsep:}\ \ignorespaces}

\newenvironment{pil}{%
  \begin{enumerate}[label=(\Alph*),leftmargin=*,itemsep=1pt,topsep=3pt]}%
  {\end{enumerate}}
\newenvironment{pilB}{%
  \par\smallskip\begin{multicols}{2}
  \begin{enumerate}[label=(\Alph*),leftmargin=*,itemsep=1pt,topsep=3pt]}%
  {\end{enumerate}\end{multicols}}

\newcommand{\gambar}[2][0.52\textwidth]{%
  \begin{center}\includegraphics[width=#1]{#2}\end{center}}

\linespread{1.05}

% ===================================================================
\begin{document}
\thispagestyle{empty}

\begin{center}
  \fbox{\parbox{0.6\textwidth}{\centering\bfseries\large
    LEMBAR SOAL DAN PEMBAHASAN}}
\end{center}
\vspace{0.5cm}

\begin{center}
  {\LARGE\bfseries FISIKA}
\end{center}
\vspace{0.5cm}

\begin{center}
  \begin{tabular}{@{\hspace{2cm}}ll@{\hspace{0.4cm}:@{\hspace{0.4cm}}}l}
    Jenjang      & & SMA / Sederajat               \\[0.3em]
    Hari/Tanggal & & \underline{\hspace{5cm}}      \\[0.3em]
    Waktu        & & \underline{\hspace{5cm}}      \\[0.3em]
    Paket Soal   & & 1                             \\
  \end{tabular}
\end{center}

\vspace{0.6cm}
\noindent\rule{\linewidth}{0.6pt}
\vspace{0.3cm}

\noindent\textbf{\underline{PETUNJUK UMUM}}
\begin{enumerate}[leftmargin=*,itemsep=4pt,topsep=4pt]
  \item Anda \textbf{tidak} diperbolehkan menggunakan sumber-sumber eksternal,
        termasuk internet, \textit{large language model} seperti ChatGPT, Claude,
        LLaMA, Gemini, Meta AI, dan sebagainya. Anda sangat dianjurkan untuk
        mencoba mengerjakan sendiri terlebih dahulu karena evaluasi ini dibuat
        untuk \textbf{mengukur} tingkat kepahaman; nilai $<70$ mengindikasikan
        \textbf{tidak lulus}, dan jika sudah melewati \textit{deadline} maka
        dianggap tidak mengerjakan yang berarti \textbf{tidak lulus}.

  \item Dokumen ini merupakan \textbf{lembar soal sekaligus pembahasan}; gunakan
        pada tahap \textbf{Review}. Di tahap ini Anda diharapkan mengerjakan ulang
        dengan disertakan penjelasan/pembelaan mengapa revisi jawaban adalah pilihan
        tersebut, dan tetap mencantumkan penjelasan pada soal yang walaupun pada
        penilaian sudah benar.

  \item Anda dianjurkan untuk mencantumkan caranya walaupun tidak termasuk
        penilaian, mengingat sifatnya adalah sebagai bahan pembelajaran.
        Mohon bertanggung jawab atas pekerjaan Anda sendiri.

  \item Terdapat tiga jenis soal: \textbf{Pilihan Ganda Biasa} (5 opsi, pilih 1);
        \textbf{Pilihan Ganda Majemuk Kompleks} (tabel \textbf{Benar}/\textbf{Salah},
        tiap pernyataan dinilai mandiri); dan \textbf{Isian Singkat}.

  \item Soal berjumlah \textbf{40 butir}: nomor 1--20 dari SSU-ITB~2024 Tipe~A,
        nomor 21--39 dari Tipe~B, dan nomor 40 adalah soal tambahan.
        Gunakan $g=10$ m/s$^2$ bila tidak disebutkan lain.

  \item Isilah detail informasi berikut.\\[0.3em]
        \textbf{Nama}\hspace{0.45cm};\enspace\underline{\hspace{8cm}}
\end{enumerate}

\vspace{0.6cm}
\noindent\textit{Dengan menandatangani kolom di bawah ini, saya menyatakan telah membaca,
memahami, dan menyetujui seluruh peraturan yang tercantum di atas.}

\vspace{0.6cm}
\noindent\fbox{\parbox{5cm}{\vspace{2.2cm}\centering\small Tanda Tangan Peserta\vspace{0.3cm}}}

\clearpage

% ===================================================================
%  TIPE A  (nomor 1--20)
% ===================================================================

\begin{soal}{1}  % PG Majemuk
Sebuah partikel bergerak sepanjang sumbu $x$. Pada $t=0$, partikel
mulai bergerak ke arah sumbu $x$ positif. Grafik kecepatan terhadap
waktu partikel tersebut ditunjukkan berikut.
\gambar[0.62\textwidth]{images/1/1-1.png}
Tentukan kebenaran setiap pernyataan berikut!

\par\smallskip
{\renewcommand{\arraystretch}{1.4}%
\begin{tabular}{@{}p{0.76\linewidth}cc@{}}
\hline
\textbf{Pernyataan} & \textbf{Benar} & \textbf{Salah}\\
\hline
(1)~Pada selang waktu 20--60\,s, jarak tempuh partikel sama dengan nol.
  & $\bigcirc$ & $\bigcirc$\\
(2)~Pada selang waktu 20--60\,s, percepatan partikel tetap (konstan).
  & $\bigcirc$ & $\bigcirc$\\
(3)~Pada selang waktu 40--60\,s, partikel bergerak diperlambat.
  & $\bigcirc$ & $\bigcirc$\\
(4)~Kelajuan partikel pada selang 10--20\,s sama dengan kelajuannya
    pada selang 60--70\,s.
  & $\bigcirc$ & $\bigcirc$\\
\hline
\end{tabular}}
\end{soal}
\begin{pembox}
\jawaban{(1)~\textbf{Salah}\quad(2)~\textbf{Benar}\quad(3)~\textbf{Salah}\quad(4)~\textbf{Benar}}
\konsep Pada grafik $v$--$t$: kemiringan $=$ percepatan; luas bertanda $=$ perpindahan; luas mutlak $=$ jarak tempuh.

\begin{enumerate}[leftmargin=*,topsep=3pt,itemsep=3pt]
\item \textbf{Salah.} Jarak tempuh $\neq0$; yang nol adalah \textit{perpindahan bersih} pada 20--60\,s, karena luas positif dan negatif sama.
\item \textbf{Benar.} Kecepatan berubah linear dari $+v$ ke $-v$, sehingga percepatan konstan.
\item \textbf{Salah.} Pada 40--60\,s, kecepatan bernilai negatif dan percepatan juga negatif (searah gerak) $\Rightarrow$ partikel \textit{dipercepat}.
\item \textbf{Benar.} Kelajuan 10--20\,s $= v$; kelajuan 60--70\,s $= v$.
\end{enumerate}
\end{pembox}

\begin{soal}{2}  % PG Biasa
Sebuah partikel bermassa 500 gram mula-mula bergerak dengan kelajuan
2,8 m/s. Partikel tersebut kemudian ditarik oleh gaya konstan $F$
sehingga dipercepat 2 m/s$^2$ selama 8 detik. Besar usaha yang
dilakukan oleh gaya $F$ selama dipercepat adalah \ldots\ J.
\begin{pilB}
  \item $72$ \item $76{,}8$ \item $80$ \item $86{,}4$ \item $96$
\end{pilB}
\end{soal}
\begin{pembox}
\jawaban{(D) $W = 86{,}4$ J}
\konsep Teorema usaha-energi: $W = \Delta K = \tfrac{1}{2}m(v^2-u^2)$.

$v = u + at = 2{,}8 + 2(8) = 18{,}8$ m/s.
\[ W = \tfrac{1}{2}(0{,}5)(18{,}8^2 - 2{,}8^2) = 0{,}25(353{,}44-7{,}84) = 0{,}25(345{,}6) = 86{,}4 \text{ J}. \]
\end{pembox}

\begin{soal}{3}  % PG Biasa
Sebuah bola dilepaskan dari atas gedung dan pada saat yang bersamaan
sebuah batu dilemparkan secara horizontal dari ketinggian yang sama.
Jika gesekan udara diabaikan, pernyataan yang paling benar adalah \ldots
\begin{pil}
  \item Bola mencapai tanah lebih dahulu karena tidak memiliki
        kecepatan horizontal.
  \item Batu mencapai tanah lebih dahulu karena memiliki tambahan
        kecepatan horizontal.
  \item Keduanya mencapai tanah pada saat yang sama.
  \item Laju akhir bola dan batu sesaat sebelum menyentuh tanah
        adalah sama.
  \item Trajektori (lintasan) bola dan batu adalah identik.
\end{pil}
\end{soal}
\begin{pembox}
\jawaban{(C) Keduanya mencapai tanah pada saat yang sama.}
\konsep Pada gerak peluru tanpa hambatan udara, gerak horizontal dan vertikal \textit{independen}.

Waktu jatuh hanya ditentukan gerak vertikal: $h = \tfrac{1}{2}gt^2$. Keduanya memiliki $v_{y,0}=0$ dan $a_y=g$ yang sama, sehingga waktu jatuh identik. Batu hanya memiliki tambahan kecepatan \textit{horizontal}, sehingga lintasannya berbeda dan laju akhirnya \textit{lebih besar} dari bola. Pilihan D dan E salah.
\end{pembox}

\begin{soal}{4}  % PG Majemuk
Pada sebuah eksperimen diperoleh data berupa jari-jari lintasan dan
percepatan sentripetal dari sebuah partikel yang bergerak melingkar
beraturan.
\begin{center}
\begin{tabular}{ccc}
\hline
No & Jari-jari (m) & $a_s$ (m/s$^2$)\\
\hline
1 & 0,2 & 45 \\
2 & 0,4 & 22,5 \\
3 & 0,6 & 15 \\
4 & 0,8 & 11,25 \\
5 & 1   & 9 \\
\hline
\end{tabular}
\end{center}
Tentukan kebenaran setiap pernyataan berikut!

\par\smallskip
{\renewcommand{\arraystretch}{1.4}%
\begin{tabular}{@{}p{0.76\linewidth}cc@{}}
\hline
\textbf{Pernyataan} & \textbf{Benar} & \textbf{Salah}\\
\hline
(1)~Hubungan $v^2 = a_s \cdot r$ berlaku pada setiap baris data.
  & $\bigcirc$ & $\bigcirc$\\
(2)~Kelajuan linear partikel adalah 3 m/s.
  & $\bigcirc$ & $\bigcirc$\\
(3)~Percepatan sentripetal berbanding lurus dengan jari-jari
    saat kelajuan konstan.
  & $\bigcirc$ & $\bigcirc$\\
(4)~Kelajuan sudut $\omega$ partikel sama pada setiap baris data.
  & $\bigcirc$ & $\bigcirc$\\
\hline
\end{tabular}}
\end{soal}
\begin{pembox}
\jawaban{(1)~\textbf{Benar}\quad(2)~\textbf{Benar}\quad(3)~\textbf{Salah}\quad(4)~\textbf{Salah}}
$v^2 = a_s r = 45(0{,}2)=9$, dst. $\Rightarrow v=3$ m/s konstan.
(3) Salah: $a_s = v^2/r$, jika $v$ konstan maka $a_s \propto 1/r$ (\textit{berbanding terbalik}).
(4) Salah: $\omega = v/r = 3/r$, berbeda di setiap baris (tergantung $r$).
\end{pembox}

\begin{soal}{5}  % PG Biasa
Sebuah benda bergerak melingkar beraturan. Pernyataan yang benar
mengenai percepatan benda adalah \ldots
\begin{pil}
  \item Percepatannya nol karena kelajuannya tetap.
  \item Percepatan selalu searah dengan kecepatan benda.
  \item Percepatan selalu menuju pusat lingkaran, besarnya tetap,
        tetapi arahnya berubah.
  \item Percepatan selalu menuju pusat lingkaran dan besarnya berubah.
  \item Percepatan menjauh dari pusat lingkaran dengan besar tetap.
\end{pil}
\end{soal}
\begin{pembox}
\jawaban{(C) Menuju pusat, besarnya tetap, arahnya berubah.}
\konsep $a_s = v^2/r = \omega^2 r$ (sentripetal).

Pada GMB, kelajuan $v$ dan $r$ tetap $\Rightarrow$ besar $a_s$ tetap. Arahnya selalu menuju pusat, tegak lurus kecepatan, sehingga arah percepatan terus berubah. Percepatan tidak nol (A salah), tidak searah kecepatan (B salah), tidak menjauh pusat (E salah).
\end{pembox}

\begin{soal}{6}  % PG Biasa
Dua buah balok bermassa $m_1=2$ kg dan $m_2=3$ kg terikat pada tali
yang melalui sebuah katrol. Balok $m_1$ berada di atas meja horizontal
kasar dengan $\mu_s=0{,}5$ dan $\mu_k=0{,}4$.
\gambar[0.42\textwidth]{images/1/1-2.png}
Besar gaya tegangan tali adalah \ldots\ N.
\begin{pilB}
  \item $8{,}0$ \item $12{,}4$ \item $14{,}0$ \item $16{,}8$ \item $22{,}4$
\end{pilB}
\end{soal}
\begin{pembox}
\jawaban{(D) $T = 16{,}8$ N}
$f_{s,\max} = 0{,}5(2)(10)=10$ N $< m_2 g = 30$ N $\Rightarrow$ sistem bergerak.
$f_k = 0{,}4(2)(10)=8$ N.
Percepatan: $a = \dfrac{m_2 g - f_k}{m_1+m_2} = \dfrac{22}{5} = 4{,}4$ m/s$^2$.
$T = m_1 a + f_k = 2(4{,}4)+8 = 16{,}8$ N.
\end{pembox}

\begin{soal}{7}  % PG Majemuk
Sebuah balok bermassa 2,5 kg meluncur turun sejauh 1 meter dari
puncak bidang miring dengan sudut $\alpha$ $\left(\tan\alpha=\tfrac{3}{4}\right)$.
Koefisien gesek kinetis $\mu_k=0{,}2$. Tentukan kebenaran setiap
pernyataan berikut!

\par\smallskip
{\renewcommand{\arraystretch}{1.4}%
\begin{tabular}{@{}p{0.76\linewidth}cc@{}}
\hline
\textbf{Pernyataan} & \textbf{Benar} & \textbf{Salah}\\
\hline
(1)~Komponen gaya gravitasi sepanjang bidang miring adalah
    $mg\sin\alpha=15$ N.
  & $\bigcirc$ & $\bigcirc$\\
(2)~Gaya gesek kinetik pada bidang miring adalah $f_k=4$ N.
  & $\bigcirc$ & $\bigcirc$\\
(3)~Usaha yang dilakukan gaya gesek bernilai positif.
  & $\bigcirc$ & $\bigcirc$\\
(4)~Usaha gaya nonkonservatif selama perpindahan 1 m adalah $-4$ J.
  & $\bigcirc$ & $\bigcirc$\\
\hline
\end{tabular}}
\end{soal}
\begin{pembox}
\jawaban{(1)~\textbf{Benar}\quad(2)~\textbf{Benar}\quad(3)~\textbf{Salah}\quad(4)~\textbf{Benar}}
$\sin\alpha=3/5$, $\cos\alpha=4/5$.
(1) $mg\sin\alpha = 2{,}5(10)(3/5)=15$ N. \checkmark
(2) $f_k = 0{,}2(2{,}5)(10)(4/5)=4$ N. \checkmark
(3) \textbf{Salah.} Gaya gesek berlawanan arah gerak $\Rightarrow$ usaha negatif.
(4) $W_{\text{gesek}} = -f_k \cdot s = -4(1)=-4$ J. \checkmark
\end{pembox}

\begin{soal}{8}  % PG Biasa
Seorang anak laki-laki dengan berat 50 kg berdiri di dalam lift.
\gambar[0.32\textwidth]{images/1/1-3.png}
Perhatikan pernyataan-pernyataan tentang gaya pada kaki anak:
\begin{enumerate}[label=\alph*.,leftmargin=*,itemsep=3pt,topsep=3pt]
  \item Jika lift diam, gaya pada kaki anak sama dengan beratnya.
  \item Jika lift dipercepat ke atas, gaya pada kaki lebih besar dari berat dan arah ke atas.
  \item Jika lift bergerak ke atas dengan kecepatan tetap, gaya pada kaki lebih kecil dari berat dan arahnya ke atas.
  \item Jika lift dipercepat ke bawah, gaya pada kaki lebih kecil dari berat dan arahnya ke bawah.
\end{enumerate}
Pernyataan yang benar adalah \ldots
\begin{pil}
  \item a, b, dan c \item a, b, c, dan d \item a saja \item a dan b \item Tidak ada
\end{pil}
\end{soal}
\begin{pembox}
\jawaban{(D) a dan b}
\konsep Hukum Newton II: $N - mg = ma$ (positif ke atas).

Diam: $a=0 \Rightarrow N=mg$. (a) Benar.
Dipercepat ke atas: $N=m(g+a)>mg$, arah $N$ ke atas. (b) Benar.
Kecepatan tetap ke atas: $a=0 \Rightarrow N=mg$, bukan lebih kecil. (c) Salah.
Dipercepat ke bawah: $N=m(g-a)<mg$, tetapi arah gaya normal \textit{tetap ke atas}, bukan ke bawah. (d) Salah.
\end{pembox}

\begin{soal}{9}  % PG Majemuk
Dua partikel identik A dan B bergerak melingkar. Grafik $a_s$ terhadap
$r$ ditunjukkan berikut.
\gambar[0.42\textwidth]{images/1/1-4.png}
Tentukan kebenaran setiap pernyataan berikut!

\par\smallskip
{\renewcommand{\arraystretch}{1.4}%
\begin{tabular}{@{}p{0.76\linewidth}cc@{}}
\hline
\textbf{Pernyataan} & \textbf{Benar} & \textbf{Salah}\\
\hline
(1)~Kemiringan garis A pada grafik $a_s$ vs $r$ menyatakan
    $\omega_A^2=1$ (rad/s)$^2$.
  & $\bigcirc$ & $\bigcirc$\\
(2)~Kemiringan garis B menyatakan $\omega_B^2=2$ (rad/s)$^2$.
  & $\bigcirc$ & $\bigcirc$\\
(3)~Kedua partikel memiliki $\omega_A=\omega_B=1$ rad/s.
  & $\bigcirc$ & $\bigcirc$\\
(4)~Selisih kelajuan sudut adalah $(\sqrt{2}-1)$ rad/s.
  & $\bigcirc$ & $\bigcirc$\\
\hline
\end{tabular}}
\end{soal}
\begin{pembox}
\jawaban{(1)~\textbf{Benar}\quad(2)~\textbf{Benar}\quad(3)~\textbf{Salah}\quad(4)~\textbf{Benar}}
\konsep $a_s = \omega^2 r$, sehingga grafik $a_s$ vs $r$ adalah garis lurus melalui origin dengan kemiringan $\omega^2$.

$\omega_A^2=1 \Rightarrow \omega_A=1$ rad/s;\; $\omega_B^2=2 \Rightarrow \omega_B=\sqrt{2}$ rad/s.
(3) Salah: $\omega_A \neq \omega_B$.
(4) Selisih $= \sqrt{2}-1$ rad/s. \checkmark
\end{pembox}

\begin{soal}{10}  % Isian
Sebuah balok yang diam bermassa 4,7 kg ditarik oleh gaya horizontal
konstan $F=10$ N di atas permukaan kasar dengan $\mu_s=0{,}6$,
$\mu_k=0{,}4$.
\gambar[0.5\textwidth]{images/1/1-5.png}
Percepatan balok tersebut adalah \ldots\ m/s$^2$.
\end{soal}
\begin{pembox}
\jawaban{$a = 0$ m/s$^2$}
Gaya gesek statik maksimum: $f_{s,\max} = \mu_s mg = 0{,}6(4{,}7)(10)=28{,}2$ N.
Karena $F = 10$ N $< f_{s,\max} = 28{,}2$ N, gaya gesek statik mampu menahan balok. Resultan gaya nol $\Rightarrow a = 0$.
\end{pembox}

\begin{soal}{11}  % PG Majemuk
Perhatikan pernyataan-pernyataan berikut tentang momentum dan gerak
pusat massa. Tentukan kebenaran setiap pernyataan!

\par\smallskip
{\renewcommand{\arraystretch}{1.4}%
\begin{tabular}{@{}p{0.76\linewidth}cc@{}}
\hline
\textbf{Pernyataan} & \textbf{Benar} & \textbf{Salah}\\
\hline
(1)~Momentum sistem bisa kekal meskipun energi mekanik tidak kekal.
  & $\bigcirc$ & $\bigcirc$\\
(2)~Syarat kekekalan momentum adalah resultan gaya luar nol.
  & $\bigcirc$ & $\bigcirc$\\
(3)~Kecepatan pusat massa berubah hanya ketika ada resultan gaya
    luar yang bekerja.
  & $\bigcirc$ & $\bigcirc$\\
(4)~Momentum total sistem selalu kekal dalam segala proses.
  & $\bigcirc$ & $\bigcirc$\\
\hline
\end{tabular}}
\end{soal}
\begin{pembox}
\jawaban{(1)~\textbf{Benar}\quad(2)~\textbf{Benar}\quad(3)~\textbf{Benar}\quad(4)~\textbf{Salah}}
(1) Pada tumbukan tidak lenting sempurna, momentum kekal, energi mekanik tidak.
(2) $\sum\vec{F}_{\text{luar}}=0$ merupakan syarat perlu dan cukup.
(3) $\sum\vec{F}_{\text{luar}}=M\vec{a}_{\text{CM}}$; jika $\sum F=0$, $\vec{v}_{\text{CM}}$ tetap.
(4) Salah: jika resultan gaya luar tidak nol (misal, ada gaya luar eksternal), momentum tidak kekal.
\end{pembox}

\begin{soal}{12}  % PG Majemuk
Dua bola identik A dan B dilepaskan dari atap gedung. Bola A
dilepaskan 1 detik sebelum bola B ($g=10$ m/s$^2$, hambatan
udara diabaikan). Tentukan kebenaran setiap pernyataan (ditinjau
saat B baru saja dilepaskan)!

\par\smallskip
{\renewcommand{\arraystretch}{1.4}%
\begin{tabular}{@{}p{0.76\linewidth}cc@{}}
\hline
\textbf{Pernyataan} & \textbf{Benar} & \textbf{Salah}\\
\hline
(1)~Kecepatan bola A saat B dilepaskan adalah 10 m/s ke bawah.
  & $\bigcirc$ & $\bigcirc$\\
(2)~Bola A sudah berada 5 m di bawah titik awal saat B dilepaskan.
  & $\bigcirc$ & $\bigcirc$\\
(3)~Setelah B dilepaskan, kecepatan relatif A terhadap B berkurang.
  & $\bigcirc$ & $\bigcirc$\\
(4)~Setelah B dilepaskan, jarak vertikal antara kedua bola selalu
    tetap 5 m.
  & $\bigcirc$ & $\bigcirc$\\
\hline
\end{tabular}}
\end{soal}
\begin{pembox}
\jawaban{(1)~\textbf{Benar}\quad(2)~\textbf{Benar}\quad(3)~\textbf{Salah}\quad(4)~\textbf{Salah}}
$v_A = g(1)=10$ m/s; $s_A=\tfrac{1}{2}g(1)^2=5$ m.
Setelah B dilepaskan: percepatan relatif $= g-g = 0$ $\Rightarrow$ kecepatan relatif \textbf{tetap} $= 10-0=10$ m/s (Salah jika "berkurang").
Jarak antara dua bola: setelah $\tau$ detik, $\Delta s = 5 + 10\tau$ (bertambah). (4) Salah.
\end{pembox}

\begin{soal}{13}  % PG Majemuk
Sebuah balok bermassa 2,5 kg meluncur dari puncak bidang miring
$\left(\tan\alpha=\tfrac{3}{4}\right)$ dengan laju awal 1 m/s.
$\mu_s=0{,}50$, $\mu_k=0{,}20$.
Tentukan kebenaran setiap pernyataan berikut!

\par\smallskip
{\renewcommand{\arraystretch}{1.4}%
\begin{tabular}{@{}p{0.76\linewidth}cc@{}}
\hline
\textbf{Pernyataan} & \textbf{Benar} & \textbf{Salah}\\
\hline
(1)~Percepatan balok sepanjang bidang miring adalah $4{,}4$ m/s$^2$.
  & $\bigcirc$ & $\bigcirc$\\
(2)~Kecepatan awal balok di puncak adalah 1 m/s.
  & $\bigcirc$ & $\bigcirc$\\
(3)~Setelah meluncur 1 m, kelajuan balok adalah 4,8 m/s.
  & $\bigcirc$ & $\bigcirc$\\
(4)~Nilai $v^2$ setelah meluncur 1 m adalah 9,8 m$^2$/s$^2$.
  & $\bigcirc$ & $\bigcirc$\\
\hline
\end{tabular}}
\end{soal}
\begin{pembox}
\jawaban{(1)~\textbf{Benar}\quad(2)~\textbf{Benar}\quad(3)~\textbf{Salah}\quad(4)~\textbf{Benar}}
$a = g\sin\alpha - \mu_k g\cos\alpha = 10(3/5)-0{,}2(10)(4/5)=6-1{,}6=4{,}4$ m/s$^2$.
$v^2 = 1^2 + 2(4{,}4)(1)=9{,}8$ $\Rightarrow$ $v=\sqrt{9{,}8}\approx3{,}13$ m/s (bukan 4,8 m/s). (3) Salah.
\end{pembox}

\begin{soal}{14}  % PG Biasa
Dua buah balok $m=7{,}6$ kg dan $M=2$ kg dihubungkan tali melalui
katrol. $\mu_s=1{,}1$; lantai licin.
\gambar[0.5\textwidth]{images/1/1-6.png}
Besar gaya horizontal maksimum $F$ agar sistem masih setimbang
adalah \ldots\ N.
\begin{pilB}
  \item $83{,}6$ \item $120{,}0$ \item $136{,}0$ \item $167{,}2$ \item $200{,}0$
\end{pilB}
\end{soal}
\begin{pembox}
\jawaban{(D) $F_{\max} = 167{,}2$ N}
Dari analisis kesetimbangan: $F = 2f$. Gesekan maksimum:
$f_{\max}=\mu_s mg=1{,}1(7{,}6)(10)=83{,}6$ N.
$F_{\max}=2(83{,}6)=167{,}2$ N.
\end{pembox}

\begin{soal}{15}  % PG Biasa
Dadang melakukan perjalanan A ke B dengan 58 km/jam, kemudian
kembali B ke A dengan 75 km/jam. Laju rata-rata seluruh perjalanan
adalah \ldots\ km/jam.
\begin{pilB}
  \item $63{,}5$ \item $65{,}0$ \item $65{,}4$ \item $66{,}5$ \item $67{,}5$
\end{pilB}
\end{soal}
\begin{pembox}
\jawaban{(C) $\bar{v} \approx 65{,}4$ km/jam}
\konsep Rata-rata harmonik: $\bar{v} = \dfrac{2v_1 v_2}{v_1+v_2}$.

\[ \bar{v}=\frac{2(58)(75)}{58+75}=\frac{8700}{133}\approx65{,}4 \text{ km/jam}. \]
\end{pembox}

\begin{soal}{16}  % PG Majemuk
Grafik simpangan gelombang berjalan pada $t=0$ yang merambat ke
kanan. Frekuensi gelombang adalah 5 Hz.
\gambar[0.55\textwidth]{images/1/1-7.png}
Tentukan kebenaran setiap pernyataan berikut!

\par\smallskip
{\renewcommand{\arraystretch}{1.4}%
\begin{tabular}{@{}p{0.76\linewidth}cc@{}}
\hline
\textbf{Pernyataan} & \textbf{Benar} & \textbf{Salah}\\
\hline
(1)~Amplitudo gelombang adalah 2 m.
  & $\bigcirc$ & $\bigcirc$\\
(2)~Panjang gelombang $\lambda = 2$ m.
  & $\bigcirc$ & $\bigcirc$\\
(3)~Bilangan gelombang $k = 2\pi$ rad/m.
  & $\bigcirc$ & $\bigcirc$\\
(4)~Persamaan simpangan gelombang adalah $y=2\cos(\pi x - 10\pi t)$.
  & $\bigcirc$ & $\bigcirc$\\
\hline
\end{tabular}}
\end{soal}
\begin{pembox}
\jawaban{(1)~\textbf{Benar}\quad(2)~\textbf{Benar}\quad(3)~\textbf{Salah}\quad(4)~\textbf{Benar}}
Dari grafik: $A=2$ m, $\lambda=2$ m.
$k = 2\pi/\lambda = 2\pi/2 = \pi$ rad/m (bukan $2\pi$). (3) Salah.
$\omega = 2\pi f = 10\pi$. Persamaan: $y=2\cos(\pi x - 10\pi t)$. (4) Benar.
\end{pembox}

\begin{soal}{17}  % Isian
Gelombang bunyi dengan frekuensi 200 Hz dimasukkan ke bagian atas
tabung vertikal berisi air. Gelombang berdiri dihasilkan pada dua
ketinggian air berturut-turut 30 cm dan 90 cm. Laju gelombang bunyi
pada kolom udara tersebut adalah \ldots\ m/s.
\end{soal}
\begin{pembox}
\jawaban{$v = 240$ m/s}
\konsep Tabung terbuka-tertutup. Dua resonansi berurutan: $\Delta L = \lambda/2$.

$\Delta L = 90 - 30 = 60$ cm $= 0{,}60$ m $\Rightarrow \lambda = 1{,}20$ m.
$v = f\lambda = 200(1{,}20) = 240$ m/s.
\end{pembox}

\begin{soal}{18}  % PG Majemuk
Sebuah gelombang berjalan:
\[
  y(x,t) = (2\,\text{m})\sin\!\left(2\pi x - 3\pi t - \frac{\pi}{3}\right).
\]
Tentukan kebenaran setiap pernyataan berikut!

\par\smallskip
{\renewcommand{\arraystretch}{1.4}%
\begin{tabular}{@{}p{0.76\linewidth}cc@{}}
\hline
\textbf{Pernyataan} & \textbf{Benar} & \textbf{Salah}\\
\hline
(1)~Gelombang merambat ke arah sumbu $+x$.
  & $\bigcirc$ & $\bigcirc$\\
(2)~Amplitudo gelombang adalah 2 m.
  & $\bigcirc$ & $\bigcirc$\\
(3)~Frekuensi sudut gelombang adalah $\omega = 2\pi$ rad/s.
  & $\bigcirc$ & $\bigcirc$\\
(4)~Laju rambat gelombang adalah 1,5 m/s.
  & $\bigcirc$ & $\bigcirc$\\
\hline
\end{tabular}}
\end{soal}
\begin{pembox}
\jawaban{(1)~\textbf{Benar}\quad(2)~\textbf{Benar}\quad(3)~\textbf{Salah}\quad(4)~\textbf{Benar}}
Bentuk $A\sin(kx-\omega t+\phi)$: merambat ke $+x$. $A=2$ m, $k=2\pi$, $\omega=3\pi$ (bukan $2\pi$). (3) Salah.
$v = \omega/k = 3\pi/2\pi = 1{,}5$ m/s. (4) Benar.
\end{pembox}

\begin{soal}{19}  % Isian
Persamaan gelombang transversal:
\[
  y(x,t) = 0{,}28\sin\!\left(\frac{\pi x}{5} + 8\pi t\right)
\]
($t$ dalam sekon, $x$ dalam meter). Laju transversal tali di
$x=1{,}5$ m saat $t=0{,}5$ s adalah \ldots\ m/s.
\end{soal}
\begin{pembox}
\jawaban{$|v_y| \approx 4{,}14$ m/s}
$v_y = \partial y/\partial t = 0{,}28(8\pi)\cos\!\left(\dfrac{\pi x}{5}+8\pi t\right)$.
Untuk $x=1{,}5$, $t=0{,}5$: argumen $= 0{,}3\pi + 4\pi = 4{,}3\pi$.
$\cos(4{,}3\pi)=\cos(0{,}3\pi)\approx0{,}5878$.
$|v_y| = 0{,}28(8\pi)(0{,}5878)\approx4{,}14$ m/s.
\end{pembox}

\begin{soal}{20}  % PG Biasa — SOAL BARU
Sebuah sumber bunyi dengan frekuensi 500 Hz bergerak mendekati
pengamat yang diam dengan kecepatan 20 m/s. Cepat rambat bunyi
340 m/s. Frekuensi bunyi yang terdeteksi oleh pengamat adalah
\ldots\ Hz.
\begin{pilB}
  \item $468{,}75$ \item $500{,}00$ \item $512{,}50$ \item $531{,}25$ \item $562{,}50$
\end{pilB}
\end{soal}
\begin{pembox}
\jawaban{(D) $f' = 531{,}25$ Hz}
\konsep Efek Doppler sumber mendekati pendengar diam: $f' = f\,\dfrac{v}{v - v_s}$.

\[ f' = 500\,\frac{340}{340-20} = 500\,\frac{340}{320} = 531{,}25 \text{ Hz}. \]
\end{pembox}

% ===================================================================
%  TIPE B  (nomor 21--40)
% ===================================================================

\begin{soal}{21}  % PG Biasa
Pelempar melempar bola dengan laju 30 m/s dan sudut elevasi $\alpha$
$\left(\tan\alpha=\tfrac{3}{4}\right)$. Penangkap berjarak 60 m.
Laju minimum penangkap untuk menangkap bola adalah \ldots\ m/s.
\begin{pilB}
  \item $5{,}50$ \item $6{,}00$ \item $6{,}50$ \item $7{,}33$ \item $8{,}00$
\end{pilB}
\end{soal}
\begin{pembox}
\jawaban{(D) $v_{\min} = \tfrac{22}{3} \approx 7{,}33$ m/s}
$\sin\alpha=3/5$, $\cos\alpha=4/5$.
$v_x = 24$ m/s, $v_y=18$ m/s.
$t = 2v_y/g = 3{,}6$ s; $R = 24(3{,}6)=86{,}4$ m.
Penangkap harus berlari $86{,}4-60=26{,}4$ m dalam 3,6 s:
$v = 26{,}4/3{,}6 = 22/3 \approx 7{,}33$ m/s.
\end{pembox}

\begin{soal}{22}  % PG Biasa
Dua balok A ($m_A=2{,}7$ kg) dan B ($m_B=5{,}9$ kg) terhubung tali
dengan panjang $\ell_{pasak-A}=8{,}2$ m dan $\ell_{A-B}=5{,}1$ m,
bergerak melingkar dengan $\omega=15$ rad/s pada meja licin.
\gambar[0.5\textwidth]{images/1/1-8.png}
Perbandingan $T_a/T_b$ adalah \ldots
\begin{pilB}
  \item $1{,}00$ \item $1{,}20$ \item $1{,}28$ \item $1{,}50$ \item $1{,}75$
\end{pilB}
\end{soal}
\begin{pembox}
\jawaban{(C) $T_a/T_b \approx 1{,}28$}
$r_A = 8{,}2$ m; $r_B = 8{,}2+5{,}1=13{,}3$ m.
$T_b = m_B \omega^2 r_B$.
$T_a = T_b + m_A \omega^2 r_A$.
\[ \frac{T_a}{T_b} = \frac{m_B r_B + m_A r_A}{m_B r_B} = \frac{5{,}9(13{,}3)+2{,}7(8{,}2)}{5{,}9(13{,}3)} = \frac{78{,}47+22{,}14}{78{,}47} \approx 1{,}28. \]
\end{pembox}

\begin{soal}{23}  % PG Biasa
Partikel 4 gram bertumbukan dengan partikel 2 gram yang diam.
Keduanya menempel setelah tumbukan. Perbandingan energi hilang
terhadap energi kinetik awal adalah \ldots
\begin{pilB}
  \item $\dfrac{1}{6}$ \item $\dfrac{1}{4}$ \item $\dfrac{1}{3}$
  \item $\dfrac{1}{2}$ \item $\dfrac{2}{3}$
\end{pilB}
\end{soal}
\begin{pembox}
\jawaban{(C) $\dfrac{1}{3}$}
\konsep Tumbukan tidak lenting sempurna. Kekekalan momentum: $v_f = \dfrac{m_1}{m_1+m_2}u$.

$\dfrac{K_f}{K_i} = \dfrac{m_1}{m_1+m_2}=\dfrac{4}{6}=\dfrac{2}{3}$.
Fraksi energi hilang $= 1-\dfrac{2}{3}=\dfrac{1}{3}$.
\end{pembox}

\begin{soal}{24}  % PG Majemuk
Pada $t=0$ partikel diam di $x=0$. Partikel bergerak dengan
percepatan seperti pada grafik.
\gambar[0.55\textwidth]{images/1/1-9.png}
Tentukan kebenaran setiap pernyataan berikut!

\par\smallskip
{\renewcommand{\arraystretch}{1.4}%
\begin{tabular}{@{}p{0.76\linewidth}cc@{}}
\hline
\textbf{Pernyataan} & \textbf{Benar} & \textbf{Salah}\\
\hline
(1)~Pada t=0--5\,s, partikel bergerak ke arah sumbu $x$ positif.
  & $\bigcirc$ & $\bigcirc$\\
(2)~Pada $t=5$ s dan $t=10$ s, partikel memiliki kecepatan sama.
  & $\bigcirc$ & $\bigcirc$\\
(3)~Pada t=5--10\,s, partikel diam (tidak bergerak).
  & $\bigcirc$ & $\bigcirc$\\
(4)~Pada $t=12$ s, partikel berada di sumbu $x$ negatif.
  & $\bigcirc$ & $\bigcirc$\\
\hline
\end{tabular}}
\end{soal}
\begin{pembox}
\jawaban{(1)~\textbf{Benar}\quad(2)~\textbf{Benar}\quad(3)~\textbf{Salah}\quad(4)~\textbf{Salah}}
Dari grafik: $a=2$ m/s$^2$ (0--5\,s); $a=0$ (5--10\,s); $a=-3$ m/s$^2$ (10--15\,s).
$v(5)=10$ m/s; $v(10)=10$ m/s (tidak berubah karena $a=0$). (2) Benar.
(3) Salah: pada 5--10\,s partikel bergerak dengan $v=10$ m/s.
(4) Salah: posisi pada $t=12$ s masih positif (perlu menghitung selengkapnya).
\end{pembox}

\begin{soal}{25}  % Isian
Partikel bergerak dari pusat koordinat ke arah $+x$: 4 m/s selama
50 s, lalu 5 m/s selama 20 s, lalu berbalik arah 3 m/s selama
30 s. Besar kecepatan rata-rata dalam 100 s adalah \ldots\ m/s.
\end{soal}
\begin{pembox}
\jawaban{$|\bar{v}| = 2{,}1$ m/s}
$\Delta x = 4(50)+5(20)-3(30)=200+100-90=210$ m.
$|\bar{v}| = |\Delta x|/\Delta t = 210/100 = 2{,}1$ m/s.
(Jika yang dimaksud kelajuan rata-rata: $s=390$ m, $\bar{v}_{laju}=3{,}9$ m/s.)
\end{pembox}

\begin{soal}{26}  % PG Biasa
Balok bermassa 5 kg ditekan pada dinding dengan gaya horizontal $F$.
Koefisien gesek $\mu=0{,}3$.
\gambar[0.5\textwidth]{images/1/1-10.png}
Gaya minimum yang diperlukan agar balok tidak jatuh adalah \ldots\ N.
\begin{pilB}
  \item $100{,}0$ \item $125{,}0$ \item $133{,}3$ \item $166{,}7$ \item $200{,}0$
\end{pilB}
\end{soal}
\begin{pembox}
\jawaban{(D) $F_{\min} \approx 166{,}7$ N}
$N=F$; syarat tidak jatuh: $f = \mu N \geq mg$.
$F \geq \dfrac{mg}{\mu} = \dfrac{5(10)}{0{,}3} = 166{,}7$ N.
\end{pembox}

\begin{soal}{27}  % PG Majemuk
Bola dijatuhkan tegak lurus dari ketinggian 50 m; setiap pantulan
menghasilkan ketinggian $\tfrac{1}{2}$ ketinggian sebelumnya.
Tentukan kebenaran setiap pernyataan berikut!

\par\smallskip
{\renewcommand{\arraystretch}{1.4}%
\begin{tabular}{@{}p{0.76\linewidth}cc@{}}
\hline
\textbf{Pernyataan} & \textbf{Benar} & \textbf{Salah}\\
\hline
(1)~Lintasan total sebelum pantulan pertama adalah 50 m.
  & $\bigcirc$ & $\bigcirc$\\
(2)~Panjang lintasan total tepat setelah pantulan ke-4 adalah
    137,5 m.
  & $\bigcirc$ & $\bigcirc$\\
(3)~Panjang lintasan total tepat setelah pantulan ke-3 sudah
    melebihi 137 m.
  & $\bigcirc$ & $\bigcirc$\\
(4)~Lintasan pertama kali melebihi 137 m tepat setelah pantulan
    ke-4.
  & $\bigcirc$ & $\bigcirc$\\
\hline
\end{tabular}}
\end{soal}
\begin{pembox}
\jawaban{(1)~\textbf{Benar}\quad(2)~\textbf{Benar}\quad(3)~\textbf{Salah}\quad(4)~\textbf{Benar}}
$S_1=50$; $S_2=50+2(25)=100$; $S_3=100+2(12{,}5)=125$; $S_4=125+2(6{,}25)=137{,}5$.
(3) Salah: $S_3=125<137$.
(4) $S_4=137{,}5>137$. Pertama kali melebihi 137 m setelah pantulan ke-4.
\end{pembox}

\begin{soal}{28}  % PG Majemuk
Balok kecil (102 gram) dilepaskan dari A, melewati lintasan lengkung
$AB$ ($\tfrac{1}{4}$ lingkaran jari-jari $R$), lalu datar BC
($\mu_k=0{,}45$), tepat berhenti di C.
\gambar[0.55\textwidth]{images/1/1-11.png}
Tentukan kebenaran setiap pernyataan berikut!

\par\smallskip
{\renewcommand{\arraystretch}{1.4}%
\begin{tabular}{@{}p{0.76\linewidth}cc@{}}
\hline
\textbf{Pernyataan} & \textbf{Benar} & \textbf{Salah}\\
\hline
(1)~Energi potensial yang hilang pada AB adalah $mgR$.
  & $\bigcirc$ & $\bigcirc$\\
(2)~Kerja gaya gesek pada BC adalah $-\mu_k mg\cdot BC$.
  & $\bigcirc$ & $\bigcirc$\\
(3)~Perbandingan $R/BC = \mu_s = 0{,}5$.
  & $\bigcirc$ & $\bigcirc$\\
(4)~Perbandingan $R/BC = \mu_k = 0{,}45$.
  & $\bigcirc$ & $\bigcirc$\\
\hline
\end{tabular}}
\end{soal}
\begin{pembox}
\jawaban{(1)~\textbf{Benar}\quad(2)~\textbf{Benar}\quad(3)~\textbf{Salah}\quad(4)~\textbf{Benar}}
Kekekalan energi: $mgR = \mu_k mg(BC) \Rightarrow R/BC = \mu_k = 0{,}45$.
(3) Salah: yang digunakan adalah $\mu_k$, bukan $\mu_s$.
\end{pembox}

\begin{soal}{29}  % PG Biasa
Ambulans (1000 Hz) bergerak menuju perempatan 1 km selama 1 menit.
Frekuensi yang dideteksi sensor di perempatan adalah \ldots\ Hz
($v_{\text{bunyi}}=340$ m/s).
\begin{pilB}
  \item $951{,}5$ \item $1000{,}0$ \item $1025{,}0$ \item $1051{,}5$ \item $1100{,}0$
\end{pilB}
\end{soal}
\begin{pembox}
\jawaban{(D) $f' \approx 1051{,}5$ Hz}
$v_s = 1000/60 = 50/3 \approx 16{,}67$ m/s.
$f' = 1000\,\dfrac{340}{340-50/3} = 1000\,\dfrac{340}{970/3} = \dfrac{1020000}{970} \approx 1051{,}5$ Hz.
\end{pembox}

\begin{soal}{30}  % PG Majemuk
Balok $m_1$ di meja licin terhubung tali tak bermassa ke massa
tergantung $m_2$ melalui katrol tak bermassa.
\gambar[0.5\textwidth]{images/1/1-12.png}
Tentukan kebenaran setiap pernyataan berikut!

\par\smallskip
{\renewcommand{\arraystretch}{1.4}%
\begin{tabular}{@{}p{0.76\linewidth}cc@{}}
\hline
\textbf{Pernyataan} & \textbf{Benar} & \textbf{Salah}\\
\hline
(1)~Percepatan sistem adalah $a=\dfrac{m_2 g}{m_1+m_2}$.
  & $\bigcirc$ & $\bigcirc$\\
(2)~Besar tegangan tali sama di kedua sisi katrol.
  & $\bigcirc$ & $\bigcirc$\\
(3)~Tegangan tali adalah $T = m_2 g$.
  & $\bigcirc$ & $\bigcirc$\\
(4)~Tegangan tali adalah $T = \dfrac{m_1 m_2 g}{m_1+m_2}$.
  & $\bigcirc$ & $\bigcirc$\\
\hline
\end{tabular}}
\end{soal}
\begin{pembox}
\jawaban{(1)~\textbf{Benar}\quad(2)~\textbf{Benar}\quad(3)~\textbf{Salah}\quad(4)~\textbf{Benar}}
$m_1$: $T=m_1 a$.\; $m_2$: $m_2 g-T=m_2 a$ $\Rightarrow$ $a=\dfrac{m_2 g}{m_1+m_2}$.
$T=m_1 a = \dfrac{m_1 m_2 g}{m_1+m_2} < m_2 g$. (3) Salah.
\end{pembox}

\begin{soal}{31}  % PG Biasa
Pesawat pengebom terbang horizontal dengan $v_p=60$ m/s pada
$h=128$ m. Bom dilepaskan tepat di atas O. Truk (tinggi 3 m) berada
100 m dari O saat bom dilepaskan. Laju truk adalah \ldots\ m/s.
\gambar[0.55\textwidth]{images/1/1-13.png}
\begin{pilB}
  \item $20$ \item $25$ \item $32$ \item $40$ \item $48$
\end{pilB}
\end{soal}
\begin{pembox}
\jawaban{(D) $v_T = 40$ m/s}
Jarak jatuh vertikal bom: $128-3=125$ m.
$t = \sqrt{2(125)/10} = 5$ s.
Bom horizontal: $x_b = 60(5)=300$ m.
Truk: $100 + v_T(5)=300 \Rightarrow v_T=40$ m/s.
\end{pembox}

\begin{soal}{32}  % PG Biasa
Balok $m=2$ kg, gaya $F=50$ N dengan $\tan\theta=3/4$ ke bawah.
$\mu_s=0{,}4$, $\mu_k=0{,}3$.
\gambar[0.55\textwidth]{images/1/1-14.png}
Besar percepatan balok adalah \ldots\ m/s$^2$.
\begin{pilB}
  \item $8{,}0$ \item $10{,}0$ \item $11{,}0$ \item $12{,}5$ \item $15{,}0$
\end{pilB}
\end{soal}
\begin{pembox}
\jawaban{(D) $a = 12{,}5$ m/s$^2$}
$F_x=40$ N, $F_y=30$ N (ke bawah). $N=mg+F_y=50$ N.
$f_{s,\max}=0{,}4(50)=20$ N $< F_x=40$ N $\Rightarrow$ bergerak.
$f_k=0{,}3(50)=15$ N.
$a = (40-15)/2 = 12{,}5$ m/s$^2$.
\end{pembox}

\begin{soal}{33}  % Isian
Mobil pemadam kecepatan 72 km/jam membunyikan sirine. Sensor di
perempatan mendeteksi perubahan frekuensi tepat sebelum dan setelah
melintas sebesar 87 Hz. Cepat rambat bunyi 340 m/s. Frekuensi
sirine adalah \ldots\ Hz.
\end{soal}
\begin{pembox}
\jawaban{$f \approx 736{,}9$ Hz}
$v_s = 20$ m/s. Selisih: $f_1-f_2 = fv\,\dfrac{2v_s}{v^2-v_s^2}=87$.
\[ f = 87\,\frac{340^2-20^2}{2(340)(20)} = 87\,\frac{115200}{13600} \approx 736{,}9 \text{ Hz}. \]
\end{pembox}

\begin{soal}{34}  % Isian
Sebuah benda bermassa 6 kg memiliki energi kinetik 200 J sesaat
sebelum menyentuh tanah (dijatuhkan dari keadaan diam, hambatan
diabaikan). Ketinggian pelepasan adalah \ldots\ m.
\end{soal}
\begin{pembox}
\jawaban{$h = \dfrac{10}{3} \approx 3{,}33$ m}
$mgh = K \Rightarrow h = \dfrac{K}{mg} = \dfrac{200}{6(10)} = \dfrac{10}{3}$ m.
\end{pembox}

\begin{soal}{35}  % PG Majemuk
Sistem A (10 kg) dan B (5 kg). Benda A ditarik ke atas sehingga
sistem bergerak ke atas dengan percepatan konstan 2 m/s$^2$.
\gambar[0.32\textwidth]{images/1/1-15.png}
Tentukan kebenaran setiap pernyataan berikut!

\par\smallskip
{\renewcommand{\arraystretch}{1.4}%
\begin{tabular}{@{}p{0.76\linewidth}cc@{}}
\hline
\textbf{Pernyataan} & \textbf{Benar} & \textbf{Salah}\\
\hline
(1)~Massa total sistem A dan B adalah 15 kg.
  & $\bigcirc$ & $\bigcirc$\\
(2)~Persamaan gerak: $T_1 - Mg = Ma$, dengan $M=15$ kg.
  & $\bigcirc$ & $\bigcirc$\\
(3)~Besar tegangan tali $T_1 = 150$ N.
  & $\bigcirc$ & $\bigcirc$\\
(4)~Besar tegangan tali $T_1 = 180$ N.
  & $\bigcirc$ & $\bigcirc$\\
\hline
\end{tabular}}
\end{soal}
\begin{pembox}
\jawaban{(1)~\textbf{Benar}\quad(2)~\textbf{Benar}\quad(3)~\textbf{Salah}\quad(4)~\textbf{Benar}}
$T_1 = M(g+a)=15(12)=180$ N (bukan 150 N; 150 N adalah jika $a=0$).
\end{pembox}

\begin{soal}{36}  % PG Majemuk
Gelombang transversal:
\[
  y(t,x) = (0{,}23\,\text{m})\sin\!\left(\frac{\pi x}{5} - 8\pi t\right).
\]
Tentukan kebenaran setiap pernyataan berikut!

\par\smallskip
{\renewcommand{\arraystretch}{1.4}%
\begin{tabular}{@{}p{0.76\linewidth}cc@{}}
\hline
\textbf{Pernyataan} & \textbf{Benar} & \textbf{Salah}\\
\hline
(1)~Gelombang merambat ke arah sumbu $+x$.
  & $\bigcirc$ & $\bigcirc$\\
(2)~Amplitudo kecepatan transversal adalah $A\omega\approx5{,}78$ m/s.
  & $\bigcirc$ & $\bigcirc$\\
(3)~Besar laju transversal di $x=1{,}5$ m, $t=0{,}5$ s lebih dari 4 m/s.
  & $\bigcirc$ & $\bigcirc$\\
(4)~Besar laju transversal di $x=1{,}5$ m, $t=0{,}5$ s adalah
    $\approx3{,}40$ m/s.
  & $\bigcirc$ & $\bigcirc$\\
\hline
\end{tabular}}
\end{soal}
\begin{pembox}
\jawaban{(1)~\textbf{Benar}\quad(2)~\textbf{Benar}\quad(3)~\textbf{Salah}\quad(4)~\textbf{Benar}}
$v_y = -0{,}23(8\pi)\cos(\pi x/5 - 8\pi t)$. Untuk $x=1{,}5$, $t=0{,}5$:
argumen $=0{,}3\pi-4\pi=-3{,}7\pi$; $\cos(-3{,}7\pi)=\cos(0{,}3\pi)\approx0{,}5878$.
$|v_y|=0{,}23(8\pi)(0{,}5878)\approx3{,}40$ m/s $<4$ m/s. (3) Salah.
\end{pembox}

\begin{soal}{37}  % PG Majemuk
Gelombang stasioner:
\[
  y = 5\cos\!\left(\frac{\pi x}{3}\right)\sin(40\pi t)
\]
($x$, $y$ dalam cm; $t$ dalam detik). Tentukan kebenaran setiap
pernyataan berikut!

\par\smallskip
{\renewcommand{\arraystretch}{1.4}%
\begin{tabular}{@{}p{0.76\linewidth}cc@{}}
\hline
\textbf{Pernyataan} & \textbf{Benar} & \textbf{Salah}\\
\hline
(1)~Gelombang stasioner ini terbentuk dari superposisi dua
    gelombang berjalan berlawanan arah.
  & $\bigcirc$ & $\bigcirc$\\
(2)~Faktor spasial persamaan ini adalah $\cos(\pi x/3)$.
  & $\bigcirc$ & $\bigcirc$\\
(3)~Jarak antara dua simpul berturut-turut adalah $\lambda/4$.
  & $\bigcirc$ & $\bigcirc$\\
(4)~Jarak antara dua simpul berturut-turut adalah 3 cm.
  & $\bigcirc$ & $\bigcirc$\\
\hline
\end{tabular}}
\end{soal}
\begin{pembox}
\jawaban{(1)~\textbf{Benar}\quad(2)~\textbf{Benar}\quad(3)~\textbf{Salah}\quad(4)~\textbf{Benar}}
Simpul: $\cos(\pi x/3)=0 \Rightarrow x = 3(2n+1)/2$. Jarak dua simpul $= 3$ cm $= \lambda/2$.
(3) Salah: jarak dua simpul = $\lambda/2$, bukan $\lambda/4$.
\end{pembox}

\begin{soal}{38}  % Isian
Grafik simpangan gelombang berjalan pada $t=0$ yang merambat ke
kanan (frekuensi 5 Hz).
\gambar[0.55\textwidth]{images/1/1-16.png}
Persamaan simpangan gelombang tersebut adalah \ldots
\end{soal}
\begin{pembox}
\jawaban{$y(x,t)=2\cos(\pi x - 10\pi t)$ (satuan SI)}
$A=2$ m, $\lambda=2$ m $\Rightarrow k=\pi$ rad/m. $f=5$ Hz $\Rightarrow \omega=10\pi$ rad/s.
Merambat ke kanan: $y=A\cos(kx-\omega t) = 2\cos(\pi x - 10\pi t)$.
\end{pembox}

\begin{soal}{39}  % PG Biasa
Senar gitar nilon panjang 90 cm, dijepit kedua ujung, total 4
simpul, laju rambat 132 m/s. Frekuensi gelombang penyusun adalah
\ldots\ Hz.
\begin{pilB}
  \item $110$ \item $165$ \item $220$ \item $330$ \item $440$
\end{pilB}
\end{soal}
\begin{pembox}
\jawaban{(C) $f = 220$ Hz}
4 simpul (termasuk ujung) $\Rightarrow$ 3 perut $\Rightarrow$ harmonik $n=3$.
$\lambda = 2L/n = 2(0{,}9)/3 = 0{,}6$ m.
$f = v/\lambda = 132/0{,}6 = 220$ Hz.
\end{pembox}

\begin{soal}{40}  % Isian — SOAL BARU
Sebuah pegas dengan konstanta pegas $k=200$ N/m dimampatkan sejauh
$x=0{,}3$ m dari posisi setimbang kemudian dilepaskan untuk
mendorong sebuah bola bermassa $m=2$ kg di atas permukaan horizontal
licin. Kelajuan bola saat pegas kembali ke posisi setimbang adalah
\ldots\ m/s.
\end{soal}
\begin{pembox}
\jawaban{$v = 3$ m/s}
\konsep Kekekalan energi: energi pegas $\to$ energi kinetik bola.

\[ \tfrac{1}{2}kx^2 = \tfrac{1}{2}mv^2 \;\Rightarrow\; v = x\sqrt{\frac{k}{m}} = 0{,}3\sqrt{\frac{200}{2}} = 0{,}3\times10 = 3 \text{ m/s}. \]
\end{pembox}

\end{document}
